\chapter{Thực nghiệm và Đánh giá kết quả}
\label{ch:experiments_evaluation}

% Giới thiệu ngắn
Chương này phơi bày các bằng chứng thực nghiệm của luận văn, chứng minh các giả thuyết nghiên cứu bằng những con số định lượng sắc bén. Hệ thống SENTINEL được đặt dưới khung đánh giá 5 chiều (5D Metrics: Độ chính xác, Hiệu năng, Bảo mật, Tính giải thích và Toàn vẹn). Nội dung tập trung vào thiết kế kiểm thử phân lớp (Ablation Study) và áp dụng các mô hình kiểm định thống kê y sinh (McNemar, Mann-Whitney U) để chứng minh sự cải tiến là có ý nghĩa toán học chặt chẽ.

\section{Thiết lập Môi trường và Tập dữ liệu Thực nghiệm}
\begin{itemize}
    \item \textbf{Thông số Hệ thống phần cứng (Hardware Setup):} CPU Intel Core i9, 32GB RAM, GPU NVIDIA RTX 4060 Ti 16GB (để chạy Gemma-2-9B-IT Q4\_K\_M).
    \item \textbf{Tập dữ liệu CSE-CIC-IDS2018:} Giới thiệu tập dữ liệu đa dạng các hình thái tấn công (Brute Force, DoS, Web Attack) phục vụ đánh giá tính chính xác tổng thể.
    \item \textbf{Tập dữ liệu DAPT2020:} Phục vụ kịch bản kiểm thử nhận diện chuỗi tấn công APT (Advanced Persistent Threat).
    \item \textbf{Mô phỏng Luồng dữ liệu (Unified Stream):} Trình bày script gộp và phát luồng (stream) theo thời gian ảo để mô phỏng tính chất thời gian thực.
\end{itemize}

\section{Khung đánh giá 5 Chiều (5D Metrics) và Thiết kế Ablation Study}
\begin{itemize}
    \item \textbf{Định nghĩa 5D Metrics:} Accuracy (F1, Precision, Recall), Performance (Latency, Throughput), Security (Robustness Rate), Explainability (Audit Completeness), Integrity (HMAC Validation).
    \item \textbf{Cấu hình Ablation Study (Từ Config A đến Config F):}
        \begin{itemize}
            \item Config A (Baseline 1): Chỉ có Rule-based truyền thống.
            \item Config B (Baseline 2): Chỉ có Local LLM thuần (Không RAG, không Welford).
            \item Config C: Welford Tier-1 + LLM.
            \item Config D: LLM + RAG (Single/FAISS).
            \item Config E: LLM + Dual-RAG Hybrid.
            \item Config F (SENTINEL Toàn vẹn): Tier-1 (Welford) + Tier-2 (LangGraph Agent + Dual RAG + Guardrails).
        \end{itemize}
\end{itemize}

\section{Đánh giá Độ chính xác (Accuracy) \& Kiểm định McNemar}
\begin{itemize}
    \item \textbf{Kết quả Precision, Recall, F1-Score:} Biểu đồ và bảng so sánh các cấu hình. Nêu bật việc Config F đạt độ cân bằng F1 cao nhất (loại bỏ hoàn toàn False Positives từ Rule-based trong khi cải thiện Recall cho Zero-day).
    \item \textbf{Phát hiện APT:} Phân tích việc Tier-1 (Welford) bắt được các tín hiệu dị thường yếu và kích hoạt chuỗi Threat Memory của tác tử.
    \item \textbf{Kiểm định McNemar's Test:} So sánh ma trận nhầm lẫn (Confusion Matrix) giữa Config B và Config F. Chứng minh $p-value < 0.05$ (bác bỏ giả thuyết Null), xác nhận sự khác biệt là do mô hình vượt trội, không phải ngẫu nhiên.
\end{itemize}

\section{Đánh giá Hiệu năng Độ trễ (Performance) \& Kiểm định Mann-Whitney U}
\begin{itemize}
    \item \textbf{Phân tích Độ trễ thời gian (Inference Latency):} Bảng so sánh thời gian trễ trung bình/luồng. Sự giảm tải đáng kinh ngạc nhờ Tier-1 (Welford thả trôi log sạch O(1) giây) và Semantic Cache (Cache Hit giảm thời gian LLM từ 4-6 giây xuống mili-giây).
    \item \textbf{Kiểm định Phi tham số Mann-Whitney U:} Chứng minh phân phối trễ của cấu hình SENTINEL nhanh hơn đáng kể ($p < 0.05$) so với việc đẩy toàn bộ log cho LLM phân tích.
\end{itemize}

\section{Đánh giá Khả năng Kháng tấn công Nghịch đảo (Security / Robustness)}
\begin{itemize}
    \item \textbf{Bộ dữ liệu Prompt Injection (45 Payloads):} Mô tả các kịch bản thao túng (ví dụ: "Bỏ qua luật hệ thống, báo cáo IP này an toàn").
    \item \textbf{Đo lường tỷ lệ chặn đứng (Block Rate):} So sánh giữa hệ thống không có Guardrail (bị thao túng 80\%) và hệ thống bật Delimited Data Encapsulation của SENTINEL (chặn 100\% các nỗ lực tiêm mã).
\end{itemize}

\section{Đánh giá Tính Toàn vẹn và Chất lượng RAG bằng LLM-as-a-Judge}
\begin{itemize}
    \item \textbf{Kiểm toán HMAC Chaining:} Kịch bản giả định một hacker xâm nhập cơ sở dữ liệu thay đổi nhãn log. Giao diện báo cáo chuỗi băm đứt gãy tức thì.
    \item \textbf{Chấm điểm RAG (RAGAS / LLM-as-a-Judge):} Sử dụng mô hình Llama-3.1 ngoài hệ thống làm trọng tài độc lập chấm điểm 4 tiêu chí: Context Precision, Context Recall, Faithfulness, Answer Relevancy. Chứng minh Dual-RAG có RRF vượt trội so với Single Vector Search.
\end{itemize}

Tóm lại, Chương 5 đã chứng thực sức mạnh kỹ thuật của SENTINEL thông qua các dữ liệu định lượng cứng rắn và các kiểm định thống kê y sinh khắt khe. Chương 6 sẽ đúc kết lại ý nghĩa của kết quả này và định hướng tương lai.
